

%%%%%%%%%%%%%%%
\begin{frame}{Programmation (Python) la base ``Messagerie''}

La base \texttt{Messagerie}, gérée par un 
SGBD (disons, MySQL)

\vfill

\figSlideWithSize{prog-python.pdf}{11cm}

\vfill

Notre utilisatrice est Athénaïs\,: elle va écrire quelques
scripts Python pour exécuter ses requêtes.
\end{frame}

%%%%%%%%%%%%%%%
\begin{frame}[fragile]{Connectons-nous}

Une application doit toujours établir une connexion
avec le serveur.

\vfill

\begin{minted}[baselinestretch=.9,
               fontsize=\small,
               framesep=2mm]{python}
connexion = pymysql.connect
              ('localhost', 
               'athénaïs', 
               'motdepasse', 
               'Messagerie',
                cursorclass=pymysql.cursors.DictCursor)
\end{minted}

\vfill

Au moins 4 paramètres: machine serveur, nom et mot de passe, nom de la base.

\vfill

Et souvent des options.

\end{frame}


%%%%%%%%%%%%%%%
\begin{frame}[fragile]{Les curseurs}

La connexion nous permet de créer des \alert{curseurs}

\begin{minted}[baselinestretch=.9,
               fontsize=\small,
               framesep=2mm]{python}
curseur = connexion.cursor()
\end{minted}
\vfill

Les curseurs permettent d'exécuter des requêtes.

\vfill

\begin{minted}[baselinestretch=.9,
               fontsize=\small,
               framesep=2mm]{python}
curseur.execute("select * from Contact")
\end{minted}

\vfill

Et de récupérer le résultat

\vfill

\begin{minted}[baselinestretch=.9,
               fontsize=\small,
               framesep=2mm]{python}
for contact in curseur.fetchall():
    print(contact['prénom'], contact['nom'])
\end{minted}
\end{frame}


%%%%%%%%%%%%%%%
\begin{frame}[fragile]{Le programme complet}

\begin{minted}[baselinestretch=.9,
               fontsize=\small,
               framesep=2mm]{python}
import pymysql
import pymysql.cursors

connexion = pymysql.connect('localhost', 'athénaïs', 
                         'motdepasse', 'Messagerie',
                         cursorclass=pymysql.cursors.DictCursor)

curseur = connexion.cursor()
curseur.execute("select * from Contact")

for contact in curseur.fetchall():
    print(contact['prénom'], contact['nom'])
\end{minted}
\vfill

\alert{Important}\,: il faudrait le rendre robuste.... À vous de jouer

\end{frame}

%%%%%%%%%%%%%%%
\begin{frame}[fragile]{Une transaction d'envoi des messages}

\begin{minted}[baselinestretch=.9,
               fontsize=\small,
               framesep=2mm]{python}
# Tous les messages non envoyés
messages = connexion.cursor()
messages.execute("select * from Message where dateEnvoi is null")
for message in messages.fetchall():
    # Ici il faudrait tamponner le message !
    
    # Ici on envoie les messages à tous les destinataires
    envois = connexion.cursor()
    envois.execute("select * from Envoi as e, Contact as c "
                   +" where e.idDestinataire=c.idContact "
                   + "and  e.idMessage = %s", message['idMessage'])
    for envoi in envois.fetchall():
        mail (envoi['email'], message['contenu')
\end{minted}

\end{frame}


%%%%%%%%%%%%%%%
\begin{frame}[fragile]{Tamponnage du message}

\begin{minted}[baselinestretch=.9,
               fontsize=\small,
               framesep=2mm]{python}
# Tous les messages non envoyés
messages = connexion.cursor()
messages.execute("select * from Message where dateEnvoi is null")
for message in messages.fetchall():
    # Marquage du message
    connexion.begin()
    maj = connexion.cursor()
    maj.execute ("Update Message set dateEnvoi='2018-12-31' "
                + "where idMessage=%s", message['idMessage'])
                
    # Ici on envoie les messages à tous les destinataires
               
    # Ici on valide
    connexion.commit()
\end{minted}

\end{frame}


%%%%%%%%%%%%
\begin{frame}{À retenir}

Transposition en Python des principes de programmation déjà
vus pendant la semaine ``Programmation''

\vfill

\begin{redbullet}
  \item Tout se passe par des curseurs 
  \item Séquence systématique exécution-boucle-fermeture (implicite)
  \item Un langage interprété rend la conversion plus simple
  
\end{redbullet}

\vfill
Prochaine question\,: mon application
est-elle à l'abri de tout problème de concurrence\,?
\end{frame}
